% The SignCollect Annotation Tool v3 — Overleaf-ready LaTeX
% Converted from the print-ready HTML technical report.
% Compiles with pdfLaTeX + BibTeX (Overleaf default). Figures are native TikZ.
\documentclass[11pt]{article}

\usepackage[T1]{fontenc}
\usepackage[utf8]{inputenc}
\usepackage{newtxtext,newtxmath}        % Times-like text + math
\usepackage[a4paper,margin=22mm]{geometry}
\usepackage{graphicx}
\usepackage{booktabs}
\usepackage{caption}
\usepackage{microtype}
\usepackage{xcolor}
\usepackage{tikz}
\usetikzlibrary{arrows.meta,positioning,fit,backgrounds,calc}
\usepackage[numbers,sort&compress]{natbib}
\usepackage[hidelinks]{hyperref}

\definecolor{accent}{HTML}{1B3A6B}
\definecolor{soft}{HTML}{EEF2F8}
\definecolor{rule}{HTML}{7A93BD}
\definecolor{seggreen}{HTML}{9EC48F}

\captionsetup{font=small,labelfont=bf}
\setlength{\parindent}{1.2em}

% ---- shared TikZ styles ----
\tikzset{
  box/.style   = {draw=accent, rounded corners=3pt, align=center, inner sep=4pt, font=\footnotesize},
  fill/.style  = {box, fill=soft},
  plain/.style = {box, fill=white},
  dash/.style  = {box, fill=white, dashed},
  flow/.style  = {-{Stealth[length=2.4mm]}, draw=accent, thick},
  lbl/.style   = {font=\scriptsize\itshape, text=gray!60!black},
}
% \sub: switch the rest of a figure node's lines to a small grey caption font.
\newcommand{\sub}{\normalfont\scriptsize\color{gray!55!black}}

\title{\vspace{-1.2em}\bfseries The SignCollect Annotation Tool v3:\\[2pt]
Automatic Sign Segmentation and Gloss Spotting in the Browser}
\author{G.\ Otterspeer\\[2pt]
\normalsize\itshape SignCollect \textperiodcentered\ Global Signbank, University of Amsterdam\\[1pt]
\normalsize Technical Report \textperiodcentered\ June 2026 \textperiodcentered\
\texttt{signcollect.nl/annotation-tool}}
\date{}

\begin{document}
\maketitle

\begin{abstract}
\noindent We describe version~3 of the SignCollect annotation tool, a browser-based editor for
multi-tier annotation of sign-language video that now closes the loop between \emph{recording}
and \emph{annotation} by embedding two neural models directly into the editing workflow. When a
user drops a clip into the tool, a \emph{sign segmentation} model based on a frozen V-JEPA~2 video
backbone and a multi-stage temporal convolutional head partitions the continuous signing stream
into individual sign segments; each resulting segment is then passed to a \emph{sign-spotting}
model that encodes the clip with a self-supervised sign representation (SignRep, built on a Hiera
vision transformer) and matches it against a 10{,}254-gloss dictionary, returning a ranked list of
candidate glosses. Both models run as warm GPU inference services reached over HTTPS, so the entire
segment-then-gloss pipeline executes in seconds without the annotator leaving the page. On the
Zin-in-NGT benchmark the segmenter attains a boundary F1 of 0.801 using RGB alone (no pose
extraction). We detail the model architectures, the system design that makes them usable from a
single-page application, and the resulting human-in-the-loop annotation workflow.

\vspace{4pt}
\noindent\textbf{Keywords:} sign language, NGT, annotation tooling, temporal segmentation, sign
spotting, self-supervised video representation, Global Signbank.
\end{abstract}

\section{Introduction}
Constructing large sign-language repositories is bottlenecked less by recording than by
\emph{annotation}: raw video must be cut into individual signs and each sign must be labelled with a
gloss linked to a lexical database. The SignCollect pipeline~\citep{signcollect} addressed the
capture side with a touchless, multi-view recording setup that lets the Sign Language of the
Netherlands (NGT) lexicon in Global Signbank be extended quickly. The annotation tool described here
addresses the downstream side, and version~3 (v3) adds machine assistance to the two most repetitive
manual steps---\emph{temporal segmentation} and \emph{gloss assignment}.

The tool remains a single self-contained web application: a multi-tier timeline editor that loads a
video (and optionally an existing ELAN \texttt{.eaf} file), lets an annotator create, move and label
segments, and exports the result as EAF. v3 keeps this editor unchanged but augments it so that, on
dropping a fresh clip, the timeline is \emph{pre-populated} with predicted sign boundaries and each
predicted segment is \emph{pre-annotated} with the ten most likely glosses. The annotator's task
shifts from authoring from scratch to verifying and correcting---a substantially faster mode of work.

This report concentrates on the two models and the engineering that makes them usable from the
browser. Section~\ref{sec:overview} gives the system architecture. Sections~\ref{sec:seg}
and~\ref{sec:spot} describe the segmentation and spotting models. Section~\ref{sec:workflow}
describes the human-in-the-loop workflow and deployment, Section~\ref{sec:results} reports benchmark
results, and Section~\ref{sec:discussion} discusses limitations.

\section{System overview}\label{sec:overview}
The browser application normalises every dropped video to 25~frames per second and H.264 at
$\le$1080p (using an in-page WebAssembly build of FFmpeg) so that the models always receive a fixed
frame rate. It then issues two kinds of HTTPS requests: a single \texttt{/segment} call carrying the
whole clip, and one \texttt{/spot} call per detected segment. Neither model runs in the browser; each
is a \emph{warm} inference service---a long-lived process that loads its model weights once and keeps
them resident on the GPU---so per-request latency is dominated by computation, not model loading.

Because the GPU host is not publicly reachable, the public web server exposes the two services
through Apache reverse-proxy paths (\texttt{/sign-spotter/} and \texttt{/sign-segmenter/}) that
forward to local ports bridged to the GPU workstation by persistent SSH reverse tunnels
(Figure~\ref{fig:arch}). Warm, this yields roughly 0.8~s per gloss-spotting call and a few seconds
for whole-clip segmentation.

\begin{figure}[t]
\centering
\begin{tikzpicture}[node distance=10mm]
  \node[fill,    text width=24mm] (br) {\textbf{Browser}\\[1pt]\sub single-page app\\\sub timeline + EAF\\\sub FFmpeg.wasm $\to$ 25\,fps};
  \node[plain,   text width=24mm, right=of br] (ap) {\textbf{Apache}\\[1pt]\sub HTTPS reverse proxy\\\sub /sign-spotter\\\sub /sign-segmenter};
  \node[dash,    text width=22mm, right=of ap] (ssh) {\textbf{SSH}\\[1pt]\sub reverse tunnels\\\sub :8000 / :8001};
  \node[plain,   text width=20mm, right=of ssh] (seg) {\textbf{Segmenter}\\\sub V-JEPA 2 + MS-TCN\\\sub POST /segment};
  \node[plain,   text width=20mm, below=4mm of seg] (sp) {\textbf{Spotter}\\\sub SignRep / Hiera\\\sub POST /spot};
  \begin{scope}[on background layer]
    \node[fill, fit=(seg)(sp), inner sep=6pt, label={[font=\footnotesize\bfseries]above:GPU host (warm)}] (gpu) {};
  \end{scope}
  \draw[flow] (br) -- node[lbl,above]{HTTPS} (ap);
  \draw[flow] (ap) -- (ssh);
  \draw[flow] (ssh) -- (gpu.west);
\end{tikzpicture}
\caption{System architecture. The browser performs all editing locally and calls two stateless HTTPS
endpoints. An Apache reverse proxy forwards them across persistent SSH tunnels to warm inference
services on a GPU workstation, which keep both models resident to minimise latency.}
\label{fig:arch}
\end{figure}

\section{Sign segmentation model}\label{sec:seg}
The segmenter answers a single question for every frame: is the signer mid-sign, and where do signs
begin and end? It is deliberately \emph{RGB-only}---it needs no hand- or body-pose extraction---which
removes a brittle and expensive preprocessing stage. The model has two parts: a large \emph{frozen}
video backbone that turns frames into rich features, and a small trained head that reads those
features as a time series and emits boundaries (Figure~\ref{fig:seg}).

\subsection{Frozen V-JEPA~2 backbone}
Features come from V-JEPA~2~\citep{vjepa2}, a ViT-L video encoder (\texttt{vjepa2-vitl-fpc64-256})
pre-trained self-supervisedly on over a million hours of general internet video---\emph{not} on sign
language. V-JEPA~2 is a Joint-Embedding Predictive Architecture: rather than reconstructing pixels,
it learns to predict the latent features of masked spatiotemporal regions, yielding representations
that capture motion and scene dynamics. We use it strictly as a frozen feature extractor.

The encoder is run on non-overlapping 16-frame chunks. Each forward pass produces a temporal grid of
patch tokens which we pool into a $3\times3$ spatial grid, giving nine region tokens (one per region)
of dimension 1024 for each of eight temporal positions in the chunk. To raise the effective temporal
resolution to one token per frame, we use a \emph{phase-dual} scheme: the backbone is run twice per
chunk, once on the frames as-is and once on the same frames shifted by one, and the two phase-shifted
token streams are interleaved. The result is a sequence of one $9\times1024$ region-token vector per
frame at 25~Hz.

\subsection{Region MS-TCN head and BIO decoding}
The trained head (\texttt{RegionMSTCN}) first collapses the nine region tokens at each frame into a
single vector using a small \emph{learned-query attention pool}, so the model can weight, say, the
dominant-hand region. The pooled per-frame sequence is then processed by a three-stage
dilated-residual temporal convolutional network in the style of MS-TCN~\citep{mstcn}. Stacked dilated
convolutions give the head a receptive field of roughly ten seconds, letting it reason about a sign
in the context of its neighbours, and a temporal smoothing (T-MSE) loss discourages jittery
frame-to-frame label flips.

The head emits, per token, one of three \emph{BIO} labels---\texttt{OUT} (no sign), \texttt{BEGIN}
(first frame of a sign), \texttt{IN} (inside a sign). The token labels are upsampled to the 25~fps
frame grid, lightly cleaned, and every maximal run of non-\texttt{OUT} frames is emitted as one
\texttt{(start, end)} segment in seconds. These segments become the timeline boxes the annotator
sees. The head was pre-trained on the German MeineDGS corpus~\citep{meinedgs} and then fine-tuned on
Zin-in-NGT, selecting the checkpoint with the best validation boundary F1.

\begin{figure}[t]
\centering
\begin{tikzpicture}[node distance=8mm]
  \node[fill, text width=18mm] (vid) {\textbf{Video}\\\textbf{@ 25\,fps}\\[1pt]\sub 16-frame chunks};
  \node[fill, text width=26mm, right=of vid] (bb) {\textbf{V-JEPA 2 (frozen)}\\\sub ViT-L video encoder\\\sub phase-dual, $3{\times}3$ regions\\\sub $\to$ 1 token/frame $\cdot$ $9{\times}1024$};
  \node[plain, text width=18mm, right=of bb] (pool) {\textbf{Attention pool}\\\sub 9 regions $\to$ 1};
  \node[fill, text width=24mm, right=of pool] (tcn) {\textbf{3-stage dilated TCN}\\\sub MS-TCN head\\\sub ${\sim}10$\,s receptive field\\\sub + T-MSE smoothing};
  \node[below=12mm of pool, font=\bfseries\footnotesize] (biolab) {Per-frame BIO labels};
  \node[below=2mm of biolab, font=\ttfamily\scriptsize] (bio)
     {\fcolorbox{gray}{gray!10}{OUT}\,\fcolorbox{seggreen!70!black}{seggreen!40}{BEG}\,%
      \fcolorbox{seggreen!70!black}{seggreen!25}{IN}\,\fcolorbox{seggreen!70!black}{seggreen!25}{IN}\,%
      \fcolorbox{gray}{gray!10}{OUT}\,\fcolorbox{seggreen!70!black}{seggreen!40}{BEG}\,%
      \fcolorbox{seggreen!70!black}{seggreen!25}{IN}};
  \node[below=8mm of bio, font=\bfseries\footnotesize] (seglab) {Segments (start, end)};
  \draw[flow] (vid) -- (bb);
  \draw[flow] (bb) -- (pool);
  \draw[flow] (pool) -- (tcn);
  \draw[flow] (tcn.south) -- (biolab.north);
  \draw[flow] (bio.south) -- (seglab.north);
\end{tikzpicture}
\caption{Segmentation pipeline. A frozen V-JEPA~2 backbone turns 25~fps frames into one
$9\times1024$ region-token vector per frame; an attention pool plus a 3-stage dilated TCN (MS-TCN)
labels each frame \texttt{OUT/BEGIN/IN}; contiguous non-\texttt{OUT} runs become timeline segments.}
\label{fig:seg}
\end{figure}

\section{Sign-spotting model}\label{sec:spot}
Given a segment, the spotter proposes which lexical sign it is. It treats spotting as
\emph{nearest-neighbour retrieval in an embedding space}~\citep{momeni} rather than as fixed-class
classification, which lets the vocabulary grow simply by adding dictionary entries---no retraining
required.

\subsection{SignRep encoder}
The segment clip is encoded by SignRep~\citep{signrep}, a self-supervised sign representation built
on a Hiera hierarchical vision transformer~\citep{hiera}. SignRep is trained as a masked autoencoder
that injects sign-specific inductive priors (skeleton-based cues) during pre-training but, crucially,
requires \emph{no} keypoints at inference time; it also uses feature regularisation and a
style-agnostic adversarial loss to make the embedding robust to signer appearance. The encoder maps a
variable-length clip to a single fixed-dimensional embedding vector.

\subsection{Dictionary matching}
Each gloss in the reference lexicon is represented by a pre-computed embedding. At query time the
segment embedding is compared by cosine similarity against the full dictionary, and the ten
highest-scoring glosses are returned with their scores (Figure~\ref{fig:spot}). The deployed
dictionary is the full-coverage set of \textbf{10{,}254 glosses}, built from the Zin-in-NGT
ground-truth signs together with the Global Signbank studio clips~\citep{signbank}, so the spotter
can in principle propose any sign in that lexicon. In the editor these ten candidates appear in a
dropdown beneath the segment; clicking one fills the segment's gloss field, and hovering a candidate
previews its reference video.

\begin{figure}[t]
\centering
\begin{tikzpicture}[node distance=9mm]
  \node[fill, text width=16mm] (clip) {\textbf{Segment}\\\sub one sign clip};
  \node[plain, text width=26mm, right=of clip] (enc) {\textbf{SignRep encoder}\\\sub Hiera ViT, self-supervised\\\sub RGB only, no keypoints};
  \node[plain, text width=14mm, right=of enc] (emb) {\textbf{Embedding}\\\sub vector};
  \node[fill, text width=26mm, above right=2mm and 6mm of emb] (dict) {\textbf{Gloss dictionary}\\\sub 10{,}254 glosses\\\sub pre-computed embeddings\\\sub Zin-in-NGT + Signbank};
  \node[plain, text width=26mm, below=10mm of dict] (cos) {\textbf{Cosine match}\\\sub nearest neighbours};
  \node[plain, text width=20mm, right=of cos] (top)
     {\textbf{Top-10}\\[2pt]\scriptsize\ttfamily 1.\,GA-MAAR-A\\2.\,PT-BHAND\\3.\,AAN-JOU\\\dots\\\itshape + scores};
  \draw[flow] (clip) -- (enc);
  \draw[flow] (enc) -- (emb);
  \draw[flow] (emb.north) to[out=60,in=180] (dict.west);
  \draw[flow] (dict) -- (cos);
  \draw[flow] (cos) -- (top);
\end{tikzpicture}
\caption{Sign-spotting pipeline. A segment clip is encoded by SignRep (a Hiera-based self-supervised
model) into one embedding, which is matched by cosine similarity against 10{,}254 pre-computed gloss
embeddings; the ten nearest glosses are returned for the annotator to choose from.}
\label{fig:spot}
\end{figure}

\section{Annotation workflow and deployment}\label{sec:workflow}
The two models are wired into the editor as a single drop-to-annotate flow:
\textbf{(1)}~The annotator drops a video; it is normalised to 25~fps in-page.
\textbf{(2)}~If the clip has no existing segments, the whole video is sent once to \texttt{/segment}
while an indeterminate progress bar is shown.
\textbf{(3)}~Returned segments populate a timeline tier.
\textbf{(4)}~Each new segment is sent to \texttt{/spot}; its dropdown first shows a spinner, then the
top-10 glosses.
\textbf{(5)}~The annotator confirms or edits boundaries and picks glosses, then exports EAF.

Segmentation runs only when a clip arrives without annotations, so dropping a video alongside an
existing EAF leaves the human's work untouched; a button re-runs segmentation on demand. The spotter
is also invoked whenever the annotator creates a segment by hand, so the assistance is available
throughout editing, not only at import.

Operationally, both services are ordinary single-process HTTP servers that load their model once and
serialise GPU calls. They were deployed on a shared CUDA environment on an NVIDIA TITAN~RTX
workstation; the public site reaches them only through the reverse-proxy and SSH-tunnel path of
Figure~\ref{fig:arch}. Health endpoints report which model and dictionary are live so operators can
confirm the correct configuration is serving.

\section{Results}\label{sec:results}
Table~\ref{tab:zin} reports segmentation quality on the Zin-in-NGT test split (345 sentences). The
V-JEPA~2 segmenter reaches a boundary F1 of 0.801 and a segment F1 of 0.740 using RGB alone,
outperforming a published pose-assisted baseline and a Hiera-backbone sibling on this in-domain NGT
data. Table~\ref{tab:dgs} summarises cross-lingual transfer to held-out German (MeineDGS) signing,
where performance is lower---as expected for a model fine-tuned on NGT---but still competitive. All
figures are reproduced from the model's own evaluation; features for the headline numbers were
extracted on GPU.

\begin{table}[t]
\centering
\caption{Sign segmentation on Zin-in-NGT (345 sentences). Higher is better. Bnd = boundary;
$\pm$4f = within 4 frames; IoU = intersection-over-union.}
\label{tab:zin}
\small
\begin{tabular}{lccccc}
\toprule
Model (RGB input) & Bnd F1 & Bnd F1 $\pm$4f & mean IoU & Seg F1 & Frame Acc \\
\midrule
\textbf{V-JEPA 2 + MS-TCN (this tool)} & \textbf{0.801} & \textbf{0.908} & \textbf{0.714} & \textbf{0.740} & \textbf{0.873} \\
Hiera + MS-TCN (sibling)              & 0.783 & 0.891 & 0.717 & 0.694 & 0.870 \\
Hands-On (published, RGB+pose, zero-shot) & 0.689 & 0.634 & 0.682 & 0.489 & 0.837 \\
\bottomrule
\end{tabular}
\end{table}

\begin{table}[t]
\centering
\caption{Cross-lingual transfer to held-out MeineDGS/DGS (boundary / segment F1).}
\label{tab:dgs}
\small
\begin{tabular}{lcc}
\toprule
Model & Boundary F1 & Segment F1 \\
\midrule
V-JEPA 2 + MS-TCN (this tool) & 0.643 & 0.481 \\
Hiera + MS-TCN (sibling)      & 0.655 & 0.502 \\
\bottomrule
\end{tabular}
\end{table}

\noindent\small The spotter is evaluated qualitatively in deployment rather than on a fixed held-out
split, since its job is to surface a short, mostly-correct candidate list for human confirmation;
warm latency is about 0.8~s per segment.\normalsize

\section{Discussion and limitations}\label{sec:discussion}
The design choice that matters most for annotation throughput is that both models are
\emph{RGB-only}: removing pose estimation eliminates a fragile dependency and lets the whole pipeline
run from the normalised clip the browser already holds. Framing both tasks around a frozen,
general-purpose video backbone (V-JEPA~2) plus a light task head, and framing spotting as dictionary
retrieval, keeps the trained components small and the vocabulary open-ended.

Limitations remain. Segmentation is fine-tuned on NGT and degrades on other sign languages
(Table~\ref{tab:dgs}). The spotter returns candidates, not decisions: rare or out-of-dictionary signs
may not appear in the top ten, and the ranked list is meant to be verified, never trusted blindly.
Latency, while interactive, is bounded by serialised GPU execution, so very long clips with many
segments are spotted progressively rather than instantly. Finally, the present report quantifies
segmentation but not end-to-end annotation time saved; a user study comparing manual versus assisted
annotation is the natural next step.

\section{Conclusion}\label{sec:conclusion}
SignCollect annotation tool v3 turns a manual sign-language editor into a human-in-the-loop system:
drop a clip, and a frozen-backbone segmenter proposes the sign boundaries while a self-supervised
spotter proposes the glosses, all from the browser. Together with the SignCollect recording
pipeline~\citep{signcollect}, this shortens the path from raw video to a glossed, Signbank-linked
EAF, supporting the rapid extension of the NGT lexicon in Global Signbank.

\paragraph{Reproducibility.} Model identifiers and code references are listed in the
bibliography~\citep{tools}. Benchmark features were extracted on GPU under mixed precision; CPU
inference may shift a few borderline boundary frames.

\bibliographystyle{plainnat}
\bibliography{references}

\end{document}
